\section{Introduction}
Chronic obstructive pulmonary disease (COPD) is a progressive chronic lung disease associated with persistent respiratory symptoms, such as breathlessness \cite{mccarthyPulmonaryRehabilitationChronic2015}. Activity related breathlessness is one of the most distressing and problematic symptoms experienced by people suffering from COPD, and is closely associated with reduced capacity for exercise and physical inactivity \cite{lewthwaiteHowAssessBreathlessness2021, odonnellExertionalDyspnoeaCOPD2016} - an outcome that is of particular concern, as it has been shown in several studies that exercise can help alleviate the progression of COPD. %Needs sources%
Pulmonary rehabilitation (PR), incorporating exercise training and other patient-tailored interventions, can improve dyspnoea (breathlessness) related quality of life and exercise capacity in people with COPD \cite{mccarthyPulmonaryRehabilitationChronic2015}.

Despite the impact of breathlessness on the quality of life of COPD patients, it is inherently subjective and thus difficult to quantifiably measure using physiological signals \cite{lewthwaiteHowAssessBreathlessness2021, odonnellExertionalDyspnoeaCOPD2016}. Despite this, in recent studies, signals such as respiratory rate, heart rate and respiratory muscle activity have demonstrated associations with patient reported dyspnoea and may enable the estimation of breathlessness severity \cite{dorssersDyspneaRespiratoryEMG2026}. Wearable monitoring and telemonitoring technologies are also increasingly being studied within PR to collect physiological and symptom data, and to support personalised feedback \cite{otto-yanezPulmonaryTelerehabilitationTelemonitoring2026}.